\documentclass{article}
\usepackage{amsmath,amssymb,hyperref}
\begin{document}
\section*{Ampleness of the maximal abelian extension of the rationals}
\subsection{Definitions and mathematical statement}
Inside a fixed algebraic closure of $\mathbb Q$, let $\mathbb Q_{\rm ab}$
be the union of the finite Galois extensions of $\mathbb Q$ whose Galois
groups are abelian.

A field $K$ is called ample if every smooth geometrically integral curve
over $K$ that has a $K$-rational point has infinitely many such points.
Here a curve is a separated scheme of finite type and dimension one;
geometric integrality means integrality after extension to an algebraic
closure of $K$.

Is $\mathbb Q_{\rm ab}$ ample? That is, for every such curve $C$ over this
field, does
\[
 C(\mathbb Q_{\rm ab})\ne\varnothing
 \quad\Longrightarrow\quad \#C(\mathbb Q_{\rm ab})=\infty
\]
hold?

\subsection{Short English statement}
Over the maximal abelian extension of the rationals, does one rational point
on a smooth geometrically integral curve guarantee infinitely many?

\subsection{Sources}
L. Bary-Soroker and A. Fehm, \emph{Open Problems in the Theory of Ample
Fields} (2011), Question I and Proposition 2.1,
\url{https://arxiv.org/abs/1106.1310}.
For later context see E. Walsberg, \emph{Topological proofs of results on
large fields}, Comptes Rendus Math\'ematique 360 (2022), 1187--1192,
introduction, \url{https://doi.org/10.5802/crmath.305}.
This is an independently written formulation of the mathematical question.
\end{document}
